\subsection{The complete pro-three unit group of the ninth-root comparison field}

\begin{statement}{A.199}{Proved}
{Two explicit conjugate-unit triples complete the ninth-root unit lattice}
Retain every embedded field and literal unit convention of
A.191, A.194, and A.197.  There is an actual global unit
\(s\in\mathcal O_M^\times\), whose fixed rational
coefficient vector is retained in the accompanying supplement,
satisfying
\[
 s^3e_3e_5u_3\,g(u_3r_2r_4)
                  =e_1e_2e_4e_8r_2r_4\,g^2(u_2).
\tag{A.199.1}
\]
Here \(g(t)=zt\), \(g(z)=z\), and \(r_2,r_4\) are
the actual A.175 anti-unit cube roots in \(E_2\); the
unit \(r\) continues to denote the separate A.197 unit.

Let \(B^{++}\) be the closed subgroup of \(\mathcal U_M\)
generated by the A.194 subgroup \(B\) and both nine-element
orbits \(g^j(r),g^j(s)\), \(0\leq j<9\).
Then the literal inclusion and quotient maps give
\[
\begin{gathered}
 B^{++}=\mathcal U_M,\qquad
 \mathcal U_M/B\simeq(\Z/3)^6,\qquad
 \mathcal U_M/B^+\simeq(\Z/3)^3,\\
 \mathcal U_M/B
   =\langle r,g(r),g^2(r),s,g(s),g^2(s)\rangle_{\F_3}.
\end{gathered}
\tag{A.199.2}
\]
The six displayed classes are a basis.  In particular the
51 actual units consisting of the old 33 A.194 units and
these two full orbits generate the entire pro-three unit
completion.  A specified subset of 27 already generates it.
Every assertion about the generated groups, quotient classes,
and residue detector is unchanged by replacing \(s\) with
either of the other two cube roots of (A.199.1).
This statement identifies the indicated completed subgroup
and quotient through their field inclusions; it does not
assign an ordinary integral index to these units.
\end{statement}

\begin{proof}
We retain the field presentation
\[
 M=\Q[t,z]/(t^9-3d^6,\ z^6+z^3+1),
 \qquad d=649,\quad \alpha=t^3/d^2.
\tag{A.199.3}
\]
The notation denotes the actual degree-54 field and its
fixed embeddings proved in A.191.  The rational coefficient
supplement defines
\[
 s=\frac1{1036253459504241}
       \sum_{i=0}^{8}\sum_{j=0}^{5}A_{6i+j}t^iz^j.
\tag{A.199.4}
\]
The retained source
\path{sections/two_inert_ninth_root_second_unit_coefficients.tex}
prints all 54 integer coefficients in unambiguous decimal
blocks.  The fixed certificate reads those blocks directly.
This retains an actual root for subsequent literal norm maps.

For the A.175 coefficient convention, the root coordinate
is the actual element
\[
 w=\frac{t^5}{d^3(\zeta-1)}
     =-\frac{t^5(\zeta+2)}{3d^3},
 \qquad \zeta=z^3,
 \qquad w^{18}=-3d^6.
\tag{A.199.5}
\]
Evaluate the literal coefficient polynomials for \(r_2,r_4\)
at this \(w\).  Substitute the A.191 polynomials for
\(e_i\), and the A.159 polynomials for \(u_i\).
Reduction by the two relations in (A.199.3) proves exactly
(A.199.1).  The certificate reconstructs the printed
integers independently, performs this multiplication with
rational coefficients, and verifies equality of all 54
coordinates.  It also verifies the equality after every
literal substitution \(t\mapsto z^kt\), \(0\leq k<9\).
The right side of (A.199.1) divided by its displayed unit
factors is a product of already certified global units.
Both \(s\) and \(s^{-1}\) therefore satisfy monic cubic
equations over algebraic integers.  Hence \(s\) and all
its displayed conjugates are actual global units.
The other two roots are \(\zeta s\) and \(\zeta^2s\),
where \(\zeta=z^3\).  This factor is a cube in \(M\)
and belongs to \(B\); also \(g(\zeta)=\zeta\).
Changing the root therefore changes neither its residue
classes modulo cubes nor the subgroup generated together
with \(B\), including all conjugates.  The root-independent
subgroup theorem and the fixed coefficient convention are
related by precisely these multiplication maps.

Order the 51 units as follows:
\[
 \text{the old 33 units of (A.194.6)};\quad
      r,g(r),\ldots,g^8(r);\quad
      s,g(s),\ldots,g^8(s).
\tag{A.199.6}
\]
For each triple \((p,a,b)\) in the following table, use
the actual residue map \((t,z)\mapsto(a,b)\).
For the \(i\)-th unit \(w_i\), define \(m_i\in\{0,1,2\}\)
by \(w_i(a,b)^{(p-1)/3}=b^{3m_i}\).
Every denominator is invertible and every unit value is
nonzero in its displayed residue field.
\[
\begin{array}{r|r|r@{\qquad}r|r|r@{\qquad}r|r|r}
 p&a&b&p&a&b&p&a&b\\\hline
307&78&287&307&78&93&307&78&53\\
307&78&168&307&78&274&307&78&46\\
307&98&287&307&98&93&307&98&53\\
307&98&168&307&131&287&307&131&93\\
307&143&287&307&143&93&307&143&53\\
307&189&287&307&193&287&523&20&468\\
523&20&410&523&20&217&523&154&468\\
523&154&410&523&156&468&523&156&410\\
523&311&468&523&311&410&613&47&467
\end{array}
\tag{A.199.7}
\]
Read rows across each line.  In the resulting 27-by-51
integer matrix, the columns
\[
\begin{split}
 1,2,3,4,5,6,7,8,9,11,12,15,16,17,19,20,23,24,25,\\
                      31,32,34,35,36,43,44,45
\end{split}
\tag{A.199.8}
\]
have exact determinant
\[
                        -16257696500\equiv1\pmod3.
\tag{A.199.9}
\]
The same full matrix gives the exact ranks
\[
 \operatorname{rank}(\text{columns }1\text{--}33)=21,
 \quad
 \operatorname{rank}(\text{columns }1\text{--}42)=24,
 \quad
 \operatorname{rank}(\text{all columns})=27.
\tag{A.199.10}
\]
These are ranks of the entire indicated column sets, not
only of the selected minor.  Every conjugate entry uses
the literal field substitution \(t\mapsto ab^j\),
\(z\mapsto b\).  The certificate independently checks
all 1377 unit values and their cubic characters.

A.194 proves that \(\mathcal U_M\) is a finitely generated
compact \(\Z_3\)-module with free rank 26 and torsion
\(\mu_9\); thus its quotient by cubes has dimension 27.
The 27 residue maps in (A.199.7) factor through that quotient.
Equation (A.199.9) proves that their combined map
\[
 \mathcal U_M/3\mathcal U_M
         \longrightarrow\F_3^{27}
\tag{A.199.11}
\]
is an isomorphism: its source has dimension 27 and the
listed units already map onto the target.
The closed subgroup generated by the selected 27 units
therefore maps onto \(\mathcal U_M/3\mathcal U_M\).
Its compact quotient is a finitely generated \(\Z_3\)-module
equal to three times itself, so Nakayama's lemma makes that
quotient zero.  This proves \(B^{++}=\mathcal U_M\),
including the torsion coordinate supplied by \(e_1=z\).

Finally A.194 gives \(3\mathcal U_M\subseteq B\).
The now injective map (A.199.11) identifies
\(B/3\mathcal U_M\) with the span of all its 33
specified generators, of exact dimension 21 by
(A.199.10).  It identifies \(B^+/3\mathcal U_M\)
with the first-orbit span, of dimension 24.  The two
literal quotient sequences consequently give
\[
 \dim_{\F_3}(\mathcal U_M/B)=27-21=6,
 \qquad
 \dim_{\F_3}(\mathcal U_M/B^+)=27-24=3.
\tag{A.199.12}
\]
The six new columns retained in the nonzero minor form
the asserted basis of the first quotient.  This proves
every map and equality in (A.199.2).

The executed certificate is
\path{code/two_inert_ninth_root_full_units_certificate.py},
with retained output and data.  The coefficient supplement,
the fixed finite determinant, and the compact Nakayama
argument establish the completed-unit assertion without
using a degree-54 class-group computation.
\end{proof}

